\section{Algorithm}
\label{sec:algorithm}

We develop an agentic system built on top of \textsc{DeXposure-FM} for two operational tasks: \emph{monitoring} and \emph{decision-making}. The system follows a \emph{forecast--then--measure} paradigm: at each decision epoch $t$, we (i) construct the latest exposure network snapshot $G_t$ and node features $X_t$, (ii) use \textsc{DeXposure-FM} to generate multi-horizon predictive distributions over future networks, (iii) compute systemic risk indicators on the predicted networks, and (iv) trigger alerts and playbook-based recommendations under explicit data-quality and uncertainty gates.

\subsection{Preliminaries}
Let $G_t=(V,E_t,W_t)$ denote the directed weighted exposure graph at time $t$, where $V$ is the set of protocols, $E_t$ is the set of directed exposure links, and $W_t=\{w_{pq,t}\}$ are edge weights. Let $X_t$ denote node-level covariates (e.g., log-TVL and derived features). For each forecast horizon $h \in \mathcal{H}$, \textsc{DeXposure-FM} outputs predictive distributions over (i) edge existence probabilities $\Pr(e_{pq,t+h}=1)$, (ii) edge weights conditional on existence, and (iii) node-level changes (e.g., log-TVL increments). We denote the model output at horizon $h$ as $\mathcal{P}_{t,h}$.

We define a deterministic \emph{predicted graph construction} operator that produces a single weighted graph $\widehat{G}_{t+h}$ from $\mathcal{P}_{t,h}$. In the minimal agent, we use the expected-weight graph:
\begin{equation}
\label{eq:expected_graph}
\widehat{w}_{pq,t+h}
~:=~
\Pr(e_{pq,t+h}=1)\cdot \mathbb{E}\!\left[w_{pq,t+h}\mid e_{pq,t+h}=1\right],
\quad
\widehat{E}_{t+h}:=\{(p,q): \Pr(e_{pq,t+h}=1)>\pi_{\min}\},
\end{equation}
yielding $\widehat{G}_{t+h}=(V,\widehat{E}_{t+h},\widehat{W}_{t+h})$.

\subsection{Step 0: Data-Health Gate}
Real-world DeFi graphs and covariates can be sparse, stale, or corrupted. We therefore introduce a deterministic data-health function $\mathrm{DataHealth}(G_t,X_t)\in[0,1]$ composed of (i) freshness checks, (ii) missingness rates, (iii) outlier/discontinuity tests, and (iv) topology sanity constraints (e.g., implausible density jumps). If the data-health score falls below a threshold $\tau_{\mathrm{data}}$, the system enters \emph{safe mode}:
\begin{equation}
\label{eq:safe_mode}
DH_t := \mathrm{DataHealth}(G_t,X_t), \qquad
\mathrm{SAFE\_MODE}_t := \mathbb{I}\{DH_t < \tau_{\mathrm{data}}\}.
\end{equation}
In safe mode, the system still produces monitoring alerts but suppresses actionable recommendations (Section~\ref{sec:decision_v1}).

\subsection{Step 1: Monitoring via Forecast--then--Measure}
\label{sec:monitoring}
For each horizon $h\in\mathcal{H}$, we forecast $\mathcal{P}_{t,h}$ with \textsc{DeXposure-FM}, construct $\widehat{G}_{t+h}$ by \eqref{eq:expected_graph}, and compute a set of monitoring metrics.

\paragraph{Systemic importance score (SIS).}
We compute a protocol-level importance score $\mathrm{SIS}_p(\widehat{G}_{t+h})$ as a weighted combination of (i) normalized PageRank, (ii) tail exposure (top-$k$ outgoing concentration), and (iii) normalized log-size (e.g., log-TVL). In the minimal system we use equal weights and $k=5$, producing a ranked watchlist $\mathrm{TopK}(\mathrm{SIS}(\widehat{G}_{t+h}))$.

\paragraph{Cross-sector spillover concentration.}
Given a sector mapping $\sigma:V\to\{1,\dots,K\}$, we aggregate predicted exposures into a sector-to-sector spillover matrix $S\in\mathbb{R}^{K\times K}$ and compute the concentration of off-diagonal spillovers via a Herfindahl--Hirschman Index (HHI):
\begin{equation}
\label{eq:spillover_hhi}
\mathrm{SpillIdx}(\widehat{G}_{t+h})
~=~
\mathrm{HHI}\!\left(\left\{S_{ij}: i\neq j\right\}\right).
\end{equation}

\paragraph{Stability indicators.}
We compute additional network statistics on $\widehat{G}_{t+h}$ such as (i) TVL concentration $\mathrm{HHI}$ over node sizes, (ii) exposure concentration $\mathrm{HHI}$ over edge weights, and (iii) network density. These serve as low-variance early-warning features.

\paragraph{Stress test losses (DebtRank-style).}
We run a standard contagion stress test on $\widehat{G}_{t+h}$ under a small set of canonical scenarios $\mathcal{S}_0$ (e.g., top-protocol shock, top-5 shock, bridge-sector shock). For each scenario $s\in\mathcal{S}_0$ we compute system-level loss $L_{t+h,s}$, contagion depth, and affected-node count.

\subsection{Alert Generation}
\label{sec:alerts}
We instantiate three alert types that are interpretable, robust, and aligned with the monitoring metrics above. Let $\mathcal{B}(m)$ denote a rolling baseline band for a metric $m$ computed on historical observations (e.g., last $L$ weeks), summarized by mean $\mu_m$ and standard deviation $\sigma_m$.

\paragraph{Alert A1: SIS watchlist shift.}
Let $W_t=\mathrm{TopK}(\mathrm{SIS}(G_t))$ and $\widehat{W}_{t+h}=\mathrm{TopK}(\mathrm{SIS}(\widehat{G}_{t+h}))$. We trigger an alert if the Jaccard distance exceeds a threshold or if a new entrant has an anomalous SIS score relative to a rolling baseline.

\paragraph{Alert A2: Spillover concentration spike.}
We trigger if the predicted spillover concentration deviates sharply from baseline, e.g.,
\begin{equation}
\label{eq:spill_alert}
\mathrm{SpillIdx}(\widehat{G}_{t+h}) - \mu_{\mathrm{SpillIdx}} > z\cdot \sigma_{\mathrm{SpillIdx}},
\end{equation}
with $z=2$ in the minimal agent.

\paragraph{Alert A3: Concentration/fragility shift with stress corroboration.}
We trigger if predicted TVL or exposure concentration shifts exceed baseline bands, and optionally corroborate with rising stress-test losses:
\begin{equation}
\label{eq:hhi_alert}
\mathrm{HHI}(\widehat{G}_{t+h}) - \mu_{\mathrm{HHI}} > z\cdot \sigma_{\mathrm{HHI}}
\quad \text{and/or} \quad
L_{t+h,s} \text{ increases materially for } s\in\mathcal{S}_0.
\end{equation}

\paragraph{Evidence and attribution.}
Each alert is accompanied by an evidence bundle including: (i) the triggering metric value and baseline band, (ii) top contributing nodes/edges/sectors (computed via marginal contribution approximations on the metric functional), (iii) horizon $h$, and (iv) a confidence label that accounts for data health and forecast uncertainty.

\subsection{Step 2: Scenario Engine (Standardized Stress Tests)}
\label{sec:scenario_engine}
The scenario engine computes \emph{risk deltas} under a small, standardized set of shocks, enabling consistent triage and decision support. For each scenario $s\in\mathcal{S}_0$ and each horizon $h\in\mathcal{H}$, we evaluate the stress-test loss on the predicted graph:
\[
L_{t+h,s} := \mathrm{StressLoss}(\widehat{G}_{t+h}, s).
\]
We rank scenarios by an impact score (minimal choice: expected loss) and retain the worst-case scenario per horizon and overall.

\subsection{Step 3: Decision Agent v1 (Playbook Recommendations)}
\label{sec:decision_v1}
The decision module maps alerts and scenario results to \emph{recommendation tickets}. In the minimal viable system, decisions are advisory and executed only by humans or external governance processes.

\paragraph{Gating.}
If $\mathrm{SAFE\_MODE}_t=1$, the agent emits only \texttt{Investigate} or \texttt{Monitor} tickets (no intervention recommendations). Otherwise, it applies deterministic playbooks.

\paragraph{Playbooks.}
We implement a small set of conservative playbooks:
(i) \textbf{P1: Escalate monitoring} for Alert A1/A2 (update watchlists, increase monitoring frequency for implicated sectors/edges);
(ii) \textbf{P2: Exposure reduction recommendation} for Alert A3 with severe stress losses (recommend reducing exposure to top contributing protocols/edges/sectors);
(iii) \textbf{P3: Contingency planning} when a standardized scenario dominates the risk profile (prepare response plans tied to measurable triggers).

\paragraph{Ticket schema.}
Each ticket contains: type (monitor/investigate/recommend), trigger alert IDs, implicated horizons, targeted protocols/edges/sectors, severity, confidence, and an evidence bundle (metric deltas, baseline bands, top contributors, and scenario losses). This enables auditability and safe human-in-the-loop deployment.

\subsection{End-to-End Procedure}
Algorithm~\ref{alg:min_agent} summarizes the minimal viable agent.

\begin{algorithm}[t]
\caption{Minimal Viable \textsc{DeXposure-FM} Agent (Monitoring + Decision v1)}
\label{alg:min_agent}
\begin{algorithmic}[1]
\Require Snapshot $(G_t,X_t)$; horizons $\mathcal{H}$; baseline bands $\mathcal{B}$; sector map $\sigma$; scenarios $\mathcal{S}_0$; thresholds $(\tau_{\mathrm{data}},\pi_{\min},z)$
\Ensure Alerts $\mathcal{A}_t$, scenario results $\mathcal{R}_t$, decision tickets $\mathcal{D}_t$
\State $DH_t \gets \mathrm{DataHealth}(G_t,X_t)$; $\mathrm{SAFE\_MODE}_t \gets \mathbb{I}\{DH_t < \tau_{\mathrm{data}}\}$
\ForAll{$h \in \mathcal{H}$}
    \State $\mathcal{P}_{t,h} \gets \textsc{DeXposure-FM}.\mathrm{Forecast}(G_t,X_t,h)$
    \State $\widehat{G}_{t+h} \gets \mathrm{BuildPredictedGraph}(\mathcal{P}_{t,h};\pi_{\min})$ \Comment{Eq.~\eqref{eq:expected_graph}}
    \State $\mathrm{SIS}_{t+h} \gets \mathrm{SIS}(\widehat{G}_{t+h})$
    \State $\mathrm{SpillIdx}_{t+h} \gets \mathrm{SpilloverHHI}(\widehat{G}_{t+h},\sigma)$ \Comment{Eq.~\eqref{eq:spillover_hhi}}
    \State $\mathrm{Stab}_{t+h} \gets \mathrm{StabilityMetrics}(\widehat{G}_{t+h})$
    \ForAll{$s \in \mathcal{S}_0$}
        \State $L_{t+h,s} \gets \mathrm{StressLoss}(\widehat{G}_{t+h}, s)$
    \EndFor
\EndFor
\State $\mathcal{A}_t \gets \mathrm{GenerateAlerts}(\{\mathrm{SIS}_{t+h},\mathrm{SpillIdx}_{t+h},\mathrm{Stab}_{t+h},L_{t+h,s}\}_{h,s}; \mathcal{B}, z, DH_t)$
\State $\mathcal{R}_t \gets \mathrm{SummarizeScenarios}(\{L_{t+h,s}\}_{h,s})$ \Comment{Rank worst-case horizons/scenarios}
\If{$\mathrm{SAFE\_MODE}_t=1$}
    \State $\mathcal{D}_t \gets \mathrm{InvestigateOnly}(\mathcal{A}_t,\mathcal{R}_t)$
\Else
    \State $\mathcal{D}_t \gets \mathrm{ApplyPlaybooks}(\mathcal{A}_t,\mathcal{R}_t)$
\EndIf
\State \Return $(\mathcal{A}_t,\mathcal{R}_t,\mathcal{D}_t)$
\end{algorithmic}
\end{algorithm}

\paragraph{Implementation note.}
Algorithm~\ref{alg:min_agent} is intentionally conservative: it requires explicit data-health gating, uses standardized scenarios, and outputs auditable tickets rather than taking autonomous actions. This design supports safe deployment while enabling measurable improvements in early warning detection and risk-aware decision support.